% Preaspiration Across Sources — Thematic Synthesis
% !TEX program = xelatex
% !TEX encoding = UTF-8 Unicode

\section{Preaspiration Across Sources}
\label{sec:preaspiration}

Preaspiration --- the occurrence of /h/ before stops (/hp/, /ht/, /hk/) --- is
one of the most debated features in SNEA historical phonology. The eight
sources in this study show dramatically different rates of preaspiration
reporting, which tells us more about the recorders than about the language.
English has no phonemic preaspiration, and English-speaking recorders
routinely failed to hear this feature. The result is a systematic gap in the
documentary record that must be filled through comparative and morphotactic
evidence.

\subsection{Cross-Source Comparison of Preaspiration Reporting}

Table~\ref{tab:preaspiration-comparison} summarizes how each source treats
preaspiration. The contrast between Edwards (who writes it consistently) and
Williams (who almost never writes it) is the most striking.

\begin{table}[h]
\centering
\caption{Preaspiration reporting across SNEA sources}
\label{tab:preaspiration-comparison}
\begin{tabular}{llll}
\toprule
\textbf{Source} & \textbf{Year} & \textbf{/h/ written?} & \textbf{Reliability} \\
\midrule
Edwards (Mahican) & 1787 & YES --- consistent & Most reliable \\
Prince-Speck (Mohegan-Pequot) & 1904 & YES --- explicit /hC/ clusters & High \\
Eliot (Wampanoag) & 1666 & Partial --- ~15\% of expected positions & Moderate \\
Winslow (Wampanoag) & 1624 & Sometimes --- \textit{kiehtan}, \textit{hobbamock} & Partial \\
Williams (Narragansett) & 1643 & Rarely --- 0 of 2,134 records & Unreliable \\
Wood (Wampanoag) & 1634 & Never & Unreliable \\
Trumbull (Wampanoag) & 1903 & Inherits Eliot's inconsistency & Moderate \\
Fielding (Mohegan) & 2013 & Not phonemic (modern loss) & N/A \\
\bottomrule
\end{tabular}
\end{table}

\subsubsection{Edwards (1787) --- The Gold Standard}

Jonathan Edwards Jr.\ is the only colonial source who can be considered
reliable for /h/ detection. He was a fluent L2 speaker of Mahican, having
learned the language as a child from playmates in Stockbridge, Massachusetts.
His transcriptions consistently write /h/ in positions where PA
reconstruction predicts it: \textit{hkeesque} 'eye', \textit{hpoon} 'winter',
\textit{keesogh} 'sun' (with /h/ in the final fricative). Edwards' training
in Hebrew and Greek, combined with his childhood fluency, gave him
phonological categories that included /h/ as a salient segment.

\subsubsection{Prince-Speck (1904) --- The Trained Linguist}

J.\ Dyneley Prince, trained in the Neogrammarian tradition, explicitly writes
/hC/ clusters where all earlier SNEA sources miss them. The spelling
\textit{bahkeder} 'bitter' shows /hk/ written; Williams or Eliot would write
only \textit{k}. Prince's Germanic training gave him an ear for aspiration
contrasts: German distinguishes aspirated /p\textsuperscript{h}, t\textsuperscript{h},
k\textsuperscript{h}/ in certain positions from unaspirated /p, t, k/ in
clusters, and the \textit{Dehnungs-h} (lengthening h) convention made him
conscious of h-like elements. Additionally, Mohegan-Pequot may have preserved
preaspiration longer than other SNEA languages, or Speck's field notes may
have captured it from Fidelia Fielding's careful speech.

\subsubsection{Eliot (1666) --- Partial and Inconsistent}

John Eliot writes \textit{hk}, \textit{hp}, \textit{ht} clusters in
approximately 15\% of expected positions. His \textit{Indian Grammar Begun}
(1666) distinguishes aspirated from unaspirated stops using English
voiced/voiceless pairs (b/p, d/t, k/g), but he does not apply this
consistently in the Bible translation. The apostrophe in Eliot's orthography
may also encode preaspiration or glottal stop --- the function is ambiguous
without comparative evidence.

\subsubsection{Winslow (1624) --- Partial but Notable}

Edward Winslow writes /h/ more consistently than any other early Wampanoag
source except Edwards. His \textit{kiehtan} (the supreme God) shows the /ht/
cluster, and \textit{hobbamock} (devil) shows word-initial /h/ where Wood's
cognate \textit{abamocho} omits it. This relatively good /h/ transcription
may reflect: (1) his intermediaries (Tisquantum, Hobbamock) had clearer /h/
articulation, (2) the Pokanoket dialect had more salient /h/ than Natick, or
(3) Winslow's ``ignorance'' of the language made him more phonetically
accurate (less normalization bias).

\subsubsection{Williams (1643) --- The Invisible /h/}

Roger Williams almost never wrote preaspirated /h/ before stops. This is the
single most important finding for the Narragansett source. In Narragansett
(and other SNEA languages), /h/ commonly precedes /p/, /t/, /k/ in
word-medial position. English speakers perceive these as single stops /k/,
/p/, /t/ because preaspiration is not a phonemic distinction in English.
Williams wrote \textit{k} where the actual pronunciation was /hk/
(preaspirated k).

The database contains only 3 records with phonetic (\textbackslash ph) fields
out of 2,134 total. All three show the same pattern:

\begin{table}[h]
\centering
\caption{The three \textbackslash ph entries in the Williams corpus}
\label{tab:williams-ph}
\begin{tabular}{llll}
\toprule
\textbf{ID} & \textbf{Williams spelling} & \textbf{\textbackslash ph (phonetic)} & \textbf{Pattern} \\
\midrule
3834 & Aaunchemókaw & ąːčimoːhkaːw & \textit{k} $\rightarrow$ /hk/ \\
4286 & Anakáusichick & anoːhkaːsiːčiːk & \textit{k} $\rightarrow$ /hk/ \\
4285 & Anakáusu & anoːhkaːsəw & \textit{k} $\rightarrow$ /hk/ \\
\bottomrule
\end{tabular}
\end{table}

All three show Williams writing \textit{k} where the reconstruction has /hk/.
The /h/ is invisible to English ears. If this pattern generalizes --- and
comparative evidence from Massachusett and Mohegan-Pequot suggests it does ---
then more than 50\% of Williams' stops may be missing a preceding /h/.

\subsection{Cowan (1973): Indirect Evidence for Preaspiration}

William Cowan's landmark study ``The Development of PA *-k- and *-hk- in
Narragansett'' (1973) demonstrated that PA *-k- (simple velar) and *-hk-
(preaspirated velar) merged phonetically in Narragansett as /hk/ medially,
but can sometimes be distinguished through the behavior of preceding vowels.
PA *-k- shortens a preceding long vowel; PA *-hk- does not. This provides an
indirect method for reconstructing preaspiration from vowel length patterns
--- even though Williams never writes the /h/.

Cowan's finding means that vowel length in Williams' orthography is not just
a vowel quality marker; it carries information about the following consonant's
preaspiration status. A long vowel before a stop suggests the stop was
simple (PA *-k-); a short vowel before a stop suggests the stop was
preaspirated (PA *-hk-). This interaction between vowel length and
preaspiration is a critical input to the conversion workflow.

\subsection{Implications for Conversion Workflows}

The pattern is clear: preaspiration was a genuine phonological feature of
SNEA languages, but English-speaking recorders routinely failed to hear it.
The recorders who wrote it most consistently (Edwards, Prince-Speck) were
either fluent L2 speakers or trained linguists with non-English phonological
categories. This has direct implications for the conversion workflows:

\begin{enumerate}
  \item \textbf{Preaspiration must be inferred, not read.} For Williams, Wood,
    and most of Eliot, preaspiration /h/ must be inferred from morphotactic
    and comparative evidence, not from the orthography.

  \item \textbf{Morphotactic position is the primary cue.} Preaspiration is
    most common before stops in the coda of stressed syllables. A
    stop preceded by a short vowel in a stressed syllable is a candidate for
    preaspiration.

  \item \textbf{Comparative evidence from Edwards is the best guide.}
    Edwards' consistent /h/ marking in Mahican provides a model for
    predicting preaspiration in sources where it is unwritten. Where Edwards
    writes /h/ and Williams omits it, the preaspiration is probable.

  \item \textbf{Cowan's vowel-length diagnostic.} For Narragansett
    specifically, vowel length patterns can serve as an indirect cue: a short
    vowel before a stop suggests preaspiration (PA *-hk-), while a long vowel
    suggests a simple stop (PA *-k-).

  \item \textbf{PA reconstructions provide the baseline.} Where PA
    reconstruction shows */h/ in a corresponding position, preaspiration is
    expected in the daughter language unless a specific sound change removed
    it.

  \item \textbf{Prince-Speck as a key witness.} The Prince-Speck glossary
    (1904) is unique among the pre-modern SNEA sources in writing /hC/
    clusters explicitly. This makes it a key witness for preaspiration in
    SNEA historical phonology, on par with Mayhew and Bourne (Wampanoag) for
    the same feature.
\end{enumerate}

Without these inference strategies, the conversion pipeline would silently
drop a phonemic /h/ from every word where it was present but unwritten ---
producing a systematic underdifferentiation that would propagate through all
downstream phonological analysis.
